\documentclass[11pt]{article}
\usepackage{amsmath,amssymb,amsthm,mathtools}
\usepackage[margin=1in]{geometry}
\usepackage{enumitem}
\usepackage{hyperref}

\title{Step 76: Signed-to-Positive Weil Matching Gate}
\author{Formed-Layer Membrane / RH Construction Notes}
\date{}

\newtheorem{theorem}{Theorem}
\newtheorem{lemma}{Lemma}
\newtheorem{proposition}{Proposition}
\newtheorem{definition}{Definition}
\newtheorem{warning}{Warning}
\newtheorem{obligation}{Obligation}
\newtheorem{corollary}{Corollary}

\newcommand{\AZ}{\mathsf A_Z}
\newcommand{\Kcmp}{\mathsf K_{\rm cmp}}
\newcommand{\Ccore}{\mathcal C^-_{\log}}
\newcommand{\VEF}{\mathcal E}
\newcommand{\qcmp}{q_{\rm cmp}}
\newcommand{\qZ}{q_Z}
\newcommand{\qW}{q_W}

\begin{document}
\maketitle

\begin{abstract}
Step 75 proposed a log-side Weil core and candidate positive feature slots for the completed carrier: prime-power, archimedean/gamma, pole/completion, and tail. Step 76 isolates the exact signed-to-positive matching gate required to turn a classical signed explicit formula into the Douglas domination
\[
  \AZ\preceq\Kcmp.
\]
The main result is not an RH proof. It is a necessary-and-sufficient operator-form criterion: a signed explicit-formula identity earns the membrane bridge only when it can be regrouped into positive carrier features and a positive closable residual form on the completed anti-invariant response space. Trace equality, finite test positivity, or signed accounting do not suffice.
\end{abstract}

\section{Setup: the bridge to be matched}
Let $Y^-$ be the completed anti-invariant response space, or a dense form core $\Ccore\subset Y^-$. Let
\[
  V_Z:Y^-\to\mathcal Z,
  \qquad
  \AZ=V_Z^*V_Z
\]
be the zero-side anti-invariant displacement readout. Let
\[
  W_{\rm cmp}:Y^-\to\mathcal E_{\rm cmp},
  \qquad
  \Kcmp=W_{\rm cmp}^*W_{\rm cmp}
\]
be the completed carrier feature map. The RH O1 bridge is
\[
  \AZ\preceq\Kcmp.
\]
By Douglas factorization, this is equivalent to the existence of a contraction
\[
  V_Z=T W_{\rm cmp},\qquad \|T\|\le1.
\]

The proposed completed feature map has four upstream-visible slots:
\[
  W_{\rm cmp}=\begin{bmatrix}
  W_{\rm pr}\\ W_\infty\\ W_{\rm pole}\\ W_{\rm tail}
  \end{bmatrix},
\]
corresponding to prime-power, archimedean/gamma, pole/completion, and support/tail records. The central issue is that the classical explicit formula is signed, while $W_{\rm cmp}^*W_{\rm cmp}$ must be positive.

\section{Signed explicit-formula matching}
\begin{definition}[Signed explicit-formula datum]
A signed explicit-formula datum on a core $\Ccore$ is a sesquilinear identity
\[
  \VEF[f,g]=\VEF_{\rm pr}[f,g]+\VEF_\infty[f,g]+\VEF_{\rm pole}[f,g]+\VEF_{\rm tail}[f,g]-\VEF_Z[f,g]=0
\]
for $f,g\in\Ccore$, together with visibility records for every term.
\end{definition}

\begin{definition}[Positive feature matching]
A signed explicit-formula datum admits a positive feature matching if there are positive forms
\[
  q_{\rm pr},\quad q_\infty,
  \quad q_{\rm pole},\quad q_{\rm tail},
  \quad q_{\rm rep}\succeq0
\]
on $\Ccore$ such that:
\begin{enumerate}[label=(\roman*)]
\item each $q_\bullet$ is represented by an upstream-visible feature map $W_\bullet$ or by an explicitly declared defect budget;
\item the completed carrier form is
\[
  \qcmp=q_{\rm pr}+q_\infty+q_{\rm pole}+q_{\rm tail}+q_{\rm rep};
\]
\item the zero-side displacement form is
\[
  \qZ[f]=\|V_Zf\|^2;
\]
\item the signed explicit-formula identity agrees with the polarization of
\[
  \qW=\qcmp-\qZ;
\]
\item $\qW\succeq0$ on the core and has a closable extension to the completed response space.
\end{enumerate}
The term $q_{\rm rep}$ is permitted only when its construction is upstream-visible and does not depend on the target zero-confinement conclusion.
\end{definition}

\begin{theorem}[Signed-to-positive matching gate]
Assume $\Ccore$ is a form core for the completed response space and the forms above are closable. A signed explicit-formula datum earns the Douglas bridge
\[
  \AZ\preceq\Kcmp
\]
if and only if it admits a positive feature matching with
\[
  \qcmp-\qZ\succeq0
\]
on the completed domain.
Equivalently, there exists a contraction $T$ such that
\[
  V_Z=T W_{\rm cmp},\qquad \|T\|\le1.
\]
\end{theorem}

\begin{proof}
If $\qcmp-\qZ\succeq0$, then for every $y$ in the completed form domain,
\[
  \|V_Zy\|^2\le\|W_{\rm cmp}y\|^2.
\]
This is precisely the Loewner inequality $V_Z^*V_Z\preceq W_{\rm cmp}^*W_{\rm cmp}$ in form sense. Douglas factorization then gives a contraction $T$ with $V_Z=T W_{\rm cmp}$.
Conversely, if such $T$ exists, then
\[
  \|V_Zy\|^2=\|TW_{\rm cmp}y\|^2\le\|W_{\rm cmp}y\|^2,
\]
so $\qcmp-\qZ\succeq0$. The feature matching condition is exactly the requirement that the signed explicit formula be represented by the polarization of this positive residual form using upstream-visible feature slots.
\end{proof}

\section{The constructive gap}
\begin{obligation}[O1 matching obligation]
To instantiate the RH bridge on the log-side core, one must construct:
\[
  W_{\rm pr},\quad W_\infty,
  \quad W_{\rm pole},\quad W_{\rm tail}
\]
and, if necessary, an upstream-positive residual feature $W_{\rm rep}$ such that
\[
  \qcmp=\|W_{\rm pr}\cdot\|^2+\|W_\infty\cdot\|^2+\|W_{\rm pole}\cdot\|^2+\|W_{\rm tail}\cdot\|^2+\|W_{\rm rep}\cdot\|^2
\]
and
\[
  \qcmp[f]-\qZ[f]
\]
is exactly the completed Weil residual form on $\Ccore$ and is nonnegative.
\end{obligation}

\begin{warning}[Weil positivity versus constructive matching]
Full Weil positivity on the exact completed anti-invariant response space is equivalent to the Douglas domination once the residual form has been identified with $\qcmp-\qZ$. However, this is not yet a constructive positive feature representation. A proof that merely states positivity of a signed form may be equivalent in strength to RH; the membrane framework demands the stronger audit object: upstream-visible positive features with bridge, tail, and null-mode records.
\end{warning}

\section{No-overread lemmas}
\begin{lemma}[Trace equality does not imply domination]
There exist positive matrices $A,K\succeq0$ with $\operatorname{tr}A=\operatorname{tr}K$ but $A\npreceq K$.
\end{lemma}
\begin{proof}
Take
\[
  A=\begin{pmatrix}2&0\\0&0\end{pmatrix},
  \qquad
  K=\begin{pmatrix}1&0\\0&1\end{pmatrix}.
\]
Both have trace $2$, but $A-K=\operatorname{diag}(1,-1)$ has a positive eigenvalue, so $A\npreceq K$.
\end{proof}

\begin{lemma}[Finite-core positivity without a core record is support-only]
Let $Y$ be infinite-dimensional and $\mathcal C_N\subset Y$ finite-dimensional. Positivity of $\qcmp-\qZ$ on each fixed $\mathcal C_N$ does not imply positivity on the completed space unless the $\mathcal C_N$ form an exhaustive form core with tail control.
\end{lemma}

\begin{proof}
Choose a vector $e$ outside the tested subspace and define a negative residual on $\operatorname{span}\{e\}$ while keeping the residual positive on the tested subspace. Every finite check passes until $e$ is included. Without an exhaustive/tail record, this is a hidden-support defect.
\end{proof}

\begin{lemma}[Tautological residual repair]
Given $A,K\succeq0$, the choice
\[
  R=(A-K)_+
\]
forces
\[
  A\preceq K+R.
\]
This is a valid mathematical budget repair but not a proof-strengthening carrier construction unless $R$ is upstream-visible and non-target-selected.
\end{lemma}

\section{Status classification}
The signed-to-positive matching gate assigns the following statuses:
\begin{center}
\begin{tabular}{ll}
\hline
Record & Status if missing\\
\hline
Positive prime feature & incomplete carrier feature\\
Positive gamma feature & incomplete completion record\\
Pole/null feature & failed null-mode legality\\
Tail/support feature & moving-ledger support-only\\
Residual positivity & O1 bridge failed\\
Form-core completion & finite-test support-only\\
Upstream visibility & smuggled repair\\
\hline
\end{tabular}
\end{center}

\section{Conclusion}
The central RH task is now a signed-to-positive matching problem: convert the signed completed explicit formula into a positive feature representation whose residual is a nonnegative closable form on the full anti-invariant response space. Only then does Weil positivity become the Douglas bridge required by the membrane framework.

\end{document}
